We propose a new method for multi-modal generative modeling that, by design, does not suffer from the limitations discussed in \Cref{sec:motivation}. Our objective is to enable both high-quality and coherent joint/conditional data generation, using a simple design (see \Cref{apdx:mld_details} for a schematic representation).
As an overview, we use deterministic uni-modal autoencoders, whereby each modality $X^i$ is encoded through its encoder $e_{\psi^i}$, which is a short form for $e^i_{\psi^i}$, into the modality specific latent variable $Z^i$ and decoded into the corresponding $\hat{X}^i=d_{\theta^i}(Z^i)$. Our approach can be interpreted as a latent variable model where the different latent variables $Z^i$ are concatenated as $Z=[Z^1,\dots,Z^M]$. This corresponds to the parametrization of the two conditional distributions as $q_\psi(z\g x)=\prod\limits_{i=1}^{M} \delta(z^i-e_{\psi^i}(x^i))$ and $p_\theta(\hat{x}\g z)=\prod\limits_{i=1}^{M}\delta(\hat{x}^i-d_{\theta^i}(z^i))$,  respectively.
Then, in place of an \gls{ELBO}, we optimize the parameters of our autoencoders by minimizing the following sum of modality specific losses: 
\begin{equation}\label{eq:AE}
    \mathcal{L}=\sum\limits_{i=1}^{M}\mathcal{L}_i,\quad \mathcal{L}_i=\int p_D^i(x^i)  l^i(x^i-d_{\theta^i}(e_{\psi^i}(x^i)))\dd x^i,
\end{equation}
where $l^i$ can be any valid distance function, e.g, the square norm $\norm{\cdot}^2$.
Parameters $\psi^i,\theta^i$ are modality specific: then, minimization of \Cref{eq:AE} corresponds to individual training of the different autoencoders.
Since the mapping from input to latent is deterministic, there is no loss of information between $X$ and $Z$.\footnote{Since the measures are not absolutely continuous w.r.t the Lebesgue measure, mutual information is $+\infty$.} Moreover, this choice avoids any form of interference in the back-propagated gradients corresponding to the uni-modal reconstruction losses. Consequently gradient conflicts issues \citep{Javaloy2022}, where stronger modalities pollute weaker ones, are avoided. 


To enable such a simple design to become a generative model, it is sufficient to generate samples from the induced latent distribution $Z\sim q_\psi(z)= \int p_D(x) q_\psi(z\g x)\dd x$ and decode them as $\hat{X}=d_\theta(Z)=[d_{\theta^1}(Z^1),\dots,d_{\theta^M}(Z^M)]$. To obtain such samples, we follow the two-stage procedure described in \cite{loaizadiagnosing,tran2021model}, where samples from the lower dimensional $q_\psi(z)$ are obtained through an appropriate generative model. We consider score-based diffusion models in latent space \citep{Rombach_2022_CVPR,Vahdat2021} to solve this task, and call our approach Multi-modal Latent Diffusion (\gls{MLD}). %Next, we delve into the details of its components.
It may be helpful to clarify, at this point, that the two-stage training of \gls{MLD} is carried out separately. Uni-modal deterministic autoencoders are pre-trained first, followed by the training of the score-based diffusion model, which is explained in more detail later.

%To conclude the overview of how our model works, for joint data generation, one can sample from noise, perform backward diffusion, and then decode the generated multi-modal latent variable to obtain the corresponding data samples. For conditional data generation, given one modality, the reverse diffusion is guided by this modality, while the other modalities are generated by sampling from noise. The generated latent variable is then decoded to obtain data samples of the missing modality.


To conclude the overview of our method, for joint data generation, one can sample from noise, perform backward diffusion, and then decode the generated multi-modal latent variable to obtain the corresponding data samples. For conditional data generation, given one modality, the reverse diffusion is guided by this modality, while the other modalities are generated by sampling from noise. The generated latent variable is then decoded to obtain data samples of the missing modality.


\subsection{Joint and Conditional Multi-modal Latent Diffusion Processes}\label{sec:mld}
In the first stage of our method, the deterministic encoders project the input modalities $X^i$ into the corresponding latent spaces $Z^i$. This transformation induces a distribution $q_\psi(z)$ for the latent variable $Z=[Z^1,\dots,Z^M]$, resulting from the concatenation of uni-modal latent variables.

\noindent \textbf{Joint generation.} To generate a new sample for all modalities we use a simple score-based diffusion model in latent space \citep{diff,Song2020,Vahdat2021,loaizadiagnosing,tran2021model}. This requires reversing a stochastic noising process, starting from a simple, Gaussian distribution. %{\color{red}initial conditions drawn from $q_\psi$}. 
Formally, the noising process is defined by a \gls{SDE} of the form:
\begin{equation}\label{eq:fw_sde_j}
    \dd R_t=\alpha(t)R_t\dd t+g(t)\dd W_t,\,\,\, R_0\sim q(r,0),
\end{equation}
where $\alpha(t)R_t$ and $g(t)$ are the drift and diffusion terms, respectively, and $W_t$ is a Wiener process. 
The time-varying probability density $q(r,t)$ of the stochastic process at time $t\in[0,T]$, where $T$ is finite, satisfies the Fokker-Planck equation \citep{oksendal2013stochastic}, with initial conditions $q(r,0)$. We assume uniqueness and existence of a stationary distribution $\rho(r)$ for the process \Cref{eq:fw_sde_j}.\footnote{This is not necessary for the validity of the method \cite{song2021maximum}} 
The forward diffusion dynamics depend on the initial conditions $R_0\sim q(r,0)$. We consider $R_0=Z$ to be the initial condition for the diffusion process, which is equivalent to $q(r,0)=q_\psi(r)$. Under loose conditions \citep{anderson1982reverse}, a time-reversed stochastic process exists, with a new \gls{SDE} of the form:
\begin{equation}\label{eq:bw_sde_j}
  \dd R_t=\left(-\alpha(T-t)R_t+g^2(T-t)\nabla\log(q(R_t,T-t))\right)\dd t+g(T-t)\dd W_t, \,\,\, R_0\sim q(r,T),  
\end{equation}
indicating that, in principle, simulation of \Cref{eq:bw_sde_j} allows to generate samples from the desired distribution $q(r,0)$. In practice, we use a \textbf{parametric score network} $s_\chi(r,t)$ to approximate the true score function, and we approximate $q(r,T)$ with the stationary distribution $\rho(r)$. Indeed, the generated data distribution $q(r,0)$ is close (in KL sense) to the true density as described by \cite{song2021maximum, howmuch}: 
\begin{equation}\label{elbo_diff_j}
\mathrm{KL}[q_\psi(r)\g\g q(r,0)]\leq \frac{1}{2}\int\limits_0^T g^2(t)\mathbb{E}[\norm{s_\chi(R_t,t)-\nabla\log q(R_t,t)}^2]\dd t+KL[q(r,T)||\rho(r)],     
\end{equation}



where the first term on the r.h.s is referred to as score-matching objective, and is the loss over which the score network is optimized, and the second is a vanishing term for $T\rightarrow\infty$. 


To conclude, joint generation of all modalities is achieved through the simulation of the reverse-time \gls{SDE} in \Cref{eq:bw_sde_j}, followed by a simple decoding procedure. Indeed, optimally trained decoders (achieving zero in \Cref{eq:AE}) can be used to transform $Z\sim q_\psi(z)$ into samples from $\int p_\theta(x\g z)q_\psi(z)\dd z=p_D(x)$.

\noindent \textbf{Conditional generation.} Given a generic partition of all modalities into non overlapping sets $A_1 \cup A_2$, where $A_2=\left(  {\{1,\dots,M\}\setminus A_1}  \right)$, conditional generation requires samples from the conditional distribution $q_\psi(z^{A_1}\g z^{A_2})$, which are based on \textit{masked} forward and backward diffusion processes. 

Given conditioning latent modalities $z^{A_2}$, we consider a modified forward diffusion process with initial conditions $R_0=\mathcal{C}(R^{A_1}_0,R^{A_2}_0)$, with $R^{A_1}_0\sim q_\psi(r^{A_1}\g z^{A_2}),R^{A_2}_0=z^{A_2}$. The composition operation $\mathcal{C}(\cdot)$ concatenates generated ($R^{A_1}$) and conditioning latents ($z^{A_2}$). As an illustration, consider $A_1 = \{1,3,5\}$, such that $X^{A_1} = \{X^1, X^3, X^5 \}$, and $A_2 = \{2,4,6\}$ such that $X^{A_2} = \{X^2, X^4, X^6 \}$. Then, $R_0=\mathcal{C}(R^{A_1}_0,R^{A_2})=\mathcal{C}(R^{A_1}_0,z^{A_2})=[ R^1_0,z^2,R^3_0,z^4,R^5_0,z^6]$.

More formally, we define the masked forward diffusion \gls{SDE}:
\begin{equation}\label{eq:cond_fw_sde_c}
    \dd R_t=m(A_1)\odot\left[ \alpha(t)R_t\dd t+g(t)\dd W_t\right],\,\,\, q(r,0)= q_\psi(r^{A_1}\g z^{A_2})\delta(r^{A_2}-z^{A_2}).
\end{equation}


The mask $m(A_1)$ contains $M$ vectors $u^i$, one per modality, and with the corresponding cardinality. If modality $j \in A_1$, then $u^j = \mathbf{1}$, otherwise $u^j = \mathbf{0}$. Then, the effect of masking is to ``freeze'' throughout the diffusion process the part of the random variable $R_t$ corresponding to the conditioning latent modalities $z^{A_2}$. We naturally associate to this modified forward process the conditional time varying density $q(r,t\g z^{A_2})=q(r^{A_1},t\g z^{A_2})\delta(r^{A_2}-z^{A_2})$.

To sample from $q_\psi(z^{A_1}\g z^{A_2})$, we derive the reverse-time dynamics of \Cref{eq:cond_fw_sde_c} as follows:
\begin{equation}\label{eq:cond_bw_sde_c}
  \dd R_t=m(A_1)\odot\left[\left(-\alpha(T-t)R_t+g^2(T-t)\nabla\log(q(R_t,T-t\g z^{A_2}))\right)\dd t+g(T-t)\dd W_t\right], 
\end{equation}
with initial conditions $R_0=\mathcal{C}(R^{A_1}_0,z^{A_2})$ and $R^{A_1}_0\sim q(r^{A_1},T\g z^{A_2})$. Then, we approximate $q(r^{A_1},T\g z^{A_2})$ by its corresponding steady state distribution $\rho(r^{A_1})$, and the true (conditional) score function $\nabla\log(q(r,t\g z^{A_2}))$ by a conditional score network $s_\chi(r^{A_1},t\g z^{A_2})$. 









